\documentclass[aspectratio=169,11pt]{beamer}
\usetheme{default}
\usecolortheme{seagull}
\setbeamertemplate{navigation symbols}{}
\setbeamertemplate{footline}[frame number]
\setbeamerfont{frametitle}{size=\large,series=\bfseries}
\setbeamerfont{title}{size=\Large,series=\bfseries}
\usepackage{graphicx}
\usepackage{listings}
\lstset{basicstyle=\ttfamily\scriptsize,breaklines=true,columns=fullflexible,keepspaces=true}

\title{Whole-MHC pangenome graph of 754 haplotypes: where it is and how to map reads to it}
\author{Robert Hoehndorf (KAUST)}
\institute{BioHackathon Japan 2026, HLA/pangenome group}
\date{15 September 2026}

\begin{document}

\begin{frame}
  \titlepage
\end{frame}

\begin{frame}{One graph holds the MHC of 754 haplotypes from five pangenome projects}
  \begin{itemize}
    \item Haplotypes: APR (UAE Arab) 106, HPRC release 2 464, JaSaPaGe (Saudi, Japanese) 38, K-PanRef (Korean) 28, CPC (Chinese) 116, GRCh38, CHM13
    \item Input: the MHC segment cut from each haplotype by our HLA workflow (GRCh38 chr6:28.4--33.6 Mb with flanks)
    \item Built with Minigraph-Cactus (Cactus 3.3.0) on the NIG supercomputer: 3 h 43 min, 32 cores
    \item Reference path \texttt{GRCh38\#0\#ctg1} = GRCh38 chr6:28,410,701--33,580,488 (graph position + 28,410,700 = chr6 coordinate)
    \item Graph: 417,896 nodes, 579,137 edges, 5.9 Mb of sequence
  \end{itemize}
\end{frame}

\begin{frame}{All files are in one team folder on the NIG supercomputer}
  \texttt{/home/asianhla/data/upload/HLA/mhc\_graph/}
  \begin{itemize}
    \item Mapping: \texttt{MHC.gbz} (graph with haplotypes, clipped), \texttt{MHC.dist}, \texttt{MHC.shortread.withzip.min}, \texttt{MHC.shortread.zipcodes} (vg giraffe indexes)
    \item Other formats: \texttt{MHC.gfa.gz}, \texttt{MHC.full.gbz} / \texttt{.full.gfa.gz} / \texttt{.full.og} (unclipped), \texttt{MHC.full.hal} (Cactus alignment)
    \item Variants: \texttt{MHC.vcf.gz} against GRCh38 (reference path ctg1)
    \item \texttt{contig\_names.tsv}: renamed contigs (\texttt{ctg<N>}) back to original names and coordinates; \texttt{README.md}: details
    \item Reads to map: \texttt{/home/asianhla/data/upload/1000G\_MHC/cram/<sample>.mhc.cram}, 2,504 1000 Genomes samples (in progress)
    \item vg: \texttt{/home/leechuck/hla/cactus/cactus-bin-v3.3.0/bin/vg} (version 1.76.1)
  \end{itemize}
\end{frame}

\begin{frame}[fragile]{Two commands map a 1000 Genomes sample to the graph (GAF or BAM)}
\begin{lstlisting}
# run in a Slurm job (sbatch -c 8 --mem=48G); this reference copy is readable from all nodes
REF=/home/leechuck/hla/ref1kg/local/GRCh38_full_analysis_set_plus_decoy_hla.fa
VG=/home/leechuck/hla/cactus/cactus-bin-v3.3.0/bin/vg
G=/home/asianhla/data/upload/HLA/mhc_graph
IDX="-Z $G/MHC.gbz -d $G/MHC.dist \
     -m $G/MHC.shortread.withzip.min -z $G/MHC.shortread.zipcodes"
CRAM=/home/asianhla/data/upload/1000G_MHC/cram/HG00096.mhc.cram

# 1. reads from the MHC CRAM
samtools collate -u -O --reference $REF $CRAM tmp \
  | samtools fastq -F 0x900 --reference $REF -1 r1.fq.gz -2 r2.fq.gz \
      -0 /dev/null -s /dev/null -n -

# 2a. align to the graph, graph coordinates (GAF)
$VG giraffe -t 8 $IDX -f r1.fq.gz -f r2.fq.gz -o gaf > HG00096.gaf

# 2b. or BAM on contig GRCh38#0#ctg1 (chr6 position = position + 28,410,700)
$VG giraffe -t 8 $IDX -f r1.fq.gz -f r2.fq.gz -o BAM | samtools sort -o HG00096.graph.bam -
\end{lstlisting}
\end{frame}

\begin{frame}{Test: 99.5\% of HG00096 MHC read pairs align in 26 seconds}
  \begin{itemize}
    \item Input: all 657,052 read pairs that bwa had placed in the MHC, chr6 alternative haplotypes or HLA contigs (1000G high-coverage CRAM of HG00096)
    \item vg giraffe, 8 threads: 26 s, 0.7 GB memory
    \item 654,059 pairs (99.5\%) aligned; 625,590 (95.2\%) with MAPQ $\ge$ 20
    \item Script: \texttt{mhc\_graph/giraffe\_test.sh}
  \end{itemize}
\end{frame}

\begin{frame}{Four points to check before using the graph}
  \begin{itemize}
    \item The graph covers only the MHC: reads from HLA-like sequence elsewhere in the genome can be forced onto it; start from reads already placed in the MHC
    \item giraffe refuses to start if \texttt{MHC.dist} is newer than the minimizer index (after copying files): \texttt{touch MHC.shortread.withzip.min MHC.shortread.zipcodes}
    \item CHM13 is a haplotype path (\texttt{chm13\#0\#ctg1\#0}), not a second reference; K-PanRef and CPC haplotypes come from their own graphs, not from assemblies; JaSaPaGe NA18952 carries one haplotype twice
    \item Workflow and job scripts: \texttt{github.com/leechuck/pangenome-bh26} (\texttt{hla/nig/mhc\_mc\_graph.sbatch}); questions: robert.hoehndorf@kaust.edu.sa
  \end{itemize}
\end{frame}

\end{document}
